% TEMPLATE DAS NOTAS EM MATEMÁTICA APLICADA (NoMA).

% Caros autores, favor não modificar o cabeçalho abaixo.

%%%%%%%%%%%%%%%%%%%%%%%%%%%%%%%%%%%%%%%%%%%%%%%%%%%%%%%%%%%%%%%%%%%%%%%%%%%%

\documentclass[11pt,twoside]{book}

\usepackage[T1]{fontenc}
\usepackage[portuguese]{babel}
\usepackage[dvips]{graphicx,color,epsfig}
\usepackage{lscape, rotating}
\usepackage{varioref,float,amssymb,amsmath}
\usepackage{float,amsfonts,amsthm}
\usepackage{changebar,longtable}
\usepackage{fancyheadings}
\usepackage{makeidx}
\usepackage{pstricks,pstricks-add,pst-math,pst-xkey}
\usepackage{hyperref}
\usepackage{graphicx}




\makeindex

\newcommand{\B}{{\tt\symbol{92}}}
\newcommand{\til}{{\tt\symbol{126}}}
\newcommand{\chap}{{\tt\symbol{94}}}
\newcommand{\agud}{{\tt\symbol{13}}}
\newcommand{\crav}{{\tt\symbol{18}}}
\newcommand{\cat}[3]{#1 \stackrel{#2}{\mathbf{\sqcap}}#3}

\rhead[]{\fancyplain{}{\rm \thepage}}
\lhead[\fancyplain{}]{\fancyplain{}{\scriptsize \rightmark}}
\cfoot{}

\setlength{\textwidth}{12.7cm}      %
\setlength{\textheight}{21.5cm}     %
\setlength{\topmargin}{0cm}         %
\setlength{\oddsidemargin}{1.61cm}  %
\setlength{\evensidemargin}{1.61cm} %

%%%%%%%%%%%%%%%%%%%%%%%%%%%%%%%%%%%%%%%%%%%%%%%%%%%%%%%%%%%%%%%%%%%%%%%%%%%%%%%%%%%%%
% Use os comandos abaixo em seu texto para os teoremas, lemas, definições, exemplos,
% propriedades, corolários, demonstrações e marcação de fim das demonstrações

%%%%%%%%%%%%%%%%%%%%%%%%%%%%%%%%%%%%%%%%%%%%%%%%%%%%%%%%%%%%%%%%%%%%%%%%%%%%%%%%%%%%%

\renewcommand{\theequation}{\thesection.\arabic{equation}}
\newtheorem{lema}{Lema}[chapter]
\newtheorem{teore}{Teorema}[chapter]
\newtheorem{defi}{Definição}[chapter]
\newtheorem{corol}{Corolário}[chapter]
\newtheorem{ex}{Exemplo}[chapter]
\newtheorem{propri}{Propriedade}[chapter]
\newcommand{\dem}{\noindent \underline{Demonstração}: $\,$}
\newcommand{\fim}{\hfill $\rule{2.0mm}{2.0mm}$ \\}    % marcação de fim das demonstrações

\renewcommand{\sectionmark}[1]{\markright{\thesection\ #1}}

% estilo Normas-ABNT (não mudar!)
\usepackage[backend=biber, style=abnt]{biblatex}
% arquivo com as referências bibliográficas do trabalho
\addbibresource{references-e-book.bib}

%%%%%%%%%%%%%%%%%%%%%%%%%%%%%%%%%%%%%%%%%%%%%%%%%%%%%%%%%%%%%%%%%%%%%%%%%%%%%%%%%%%%%
%%%%%%%%%%%%%%%%%%%%%%%%%%%%%%%%%%%%%%%%%%%%%%%%%%%%%%%%%%%%%%%%%%%%%%%%%%%%%%%%%%%%%


\begin{document}

\pagestyle{myheadings}

\input{primeiras-paginas-e-book}


%%%%%%%%%%%%%%%%%%%%%%%%%%%%%%%%%%%%%%%%%

% Quinta Folha - Sumário

%%%%%%%%%%%%%%%%%%%%%%%%%%%%%%%%%%%%%%%%%

\tableofcontents


%%%%%%%%%%%%%%%%%%%%%%%%%%%%%%%%%%%%%%%%%%%%%%%%%%%%%%%%%%%%%%%%%%%%%%%%%%%%%%%%%%%%%%%%%%%%%%%%%%%%%%

% Verso da Última Folha do Sumário - Continuação do Sumário ou em branco

%%%%%%%%%%%%%%%%%%%%%%%%%%%%%%%%%%%%%%%%%%%%%%%%%%%%%%%%%%%%%%%%%%%%%%%%%%%%%%%%%%%%%%%%%%%%%%%%%%%%%%%

\newpage  \thispagestyle{empty}  $\frac{}{}$ % retirar caso o sumário continue no verso da folha

%%%%%%%%%%%%%%%%%%%%%%%%%%%%%%%%%%%%%%%%%%%%%%%%%%%%%%%%%%%%%%%%%%%%%%%%%%%%%%%%%%%%%%%%%%%%%%%%%%%%%%%

%%%%%%%%%%%%%%%%%%%%%%%%%%%%%%%%%%%%%%%%%%%%%%%%%%%%%%%%

% Sexta Folha - Prefácio dos autores - Obrigatória

%%%%%%%%%%%%%%%%%%%%%%%%%%%%%%%%%%%%%%%%%%%%%%%%%%%%%%%%%


\input{prefacio-e-book}  %           prefácio dos autores

%%%%%%%%%%%%%%%%%%%%%%%%%%%%%%%%%%%%%%%%%%%%%%%%%%%%%%%%%%%%%%%%%%%%%%%%%%%%%%%%%%%%%%%%%%%%%%%%%%%%%%%%%

% Verso da Última Folha do Prefácio- Continuação do Prefácio dos autores ou em branco

%%%%%%%%%%%%%%%%%%%%%%%%%%%%%%%%%%%%%%%%%%%%%%%%%%%%%%%%%%%%%%%%%%%%%%%%%%%%%%%%%%%%%%%%%%%%%%%%%%%%%%%%%

\newpage  \thispagestyle{empty}   $\frac{}{}$ % retirar caso o prefácio continue no verso da  folha

%%%%%%%%%%%%%%%%%%%%%%%%%%%%%%%%%%%%%%%%%%%%%%%%%%%%%%%%%%%%%%%%%%%%%%%%%%%%%%%%%%%%%%%%%%%%%%%%%%%%%%%%%%

\newpage

%\pagestyle{fancyplain}  %

%%%%%%%%%%%%%%%%%%%%%%%%%%%%%%%%%%%%%%%%%%%%%%%%%%%%%%%%%%%%%%%%%%%
% OBS: OS TEXTOS DEVEM TER NO MÍNIMO 80 PÁGINAS E NO MÁXIMO 150.
%%%%%%%%%%%%%%%%%%%%%%%%%%%%%%%%%%%%%%%%%%%%%%%%%%%%%%%%%%%%%%%%%%%%%%%%%%%%%%%%%%%%%%%%%%%

% O texto pode ser feito em um só arquivo ou pode ser separado por capítulos.
% Recomenda-se que seja separado por capítulos.
% Para cada capítulo  colocar o comando \input{nome do arquivo}. Exemplo:

%%%%%%%%%%%%%%%%%%%%%%%%%%%%%%%%%%%%%%%%%%%%%%%%%%%%%%%%%%%%%%%%%%%%%%%%%%%%%%%%%%%%%%%%%%%%%

\input{cap1-e-book}
\input{cap2-e-book}
\input{cap3-e-book}
%\bibliographystyle{abnt-alf}
%\bibliography{references-e-book}

\printbibliography

\printindex

%%%%%%%%%%%%%%%%%%%%%%%%%%%%%%%%%%%%%%%%%%%%%%%%%%%%%%%%%%%%%%%%%%%%%%%%%%%%%%%%
% Para fazer o índice remissivo, o preâmbulo deve conter os comandos \usepackage{makeidx}
% e \makeindex (como indicado neste exemplo). Além disso, após as referências, incluir
% o comando  \printindex.
% Para que as expressões (à escolha dos autores) apareçam no índice remissivo, colocar o
% comando \index( expressão desejada) próximo ao local onde o assunto é tratado.
% O latex fará, automaticamente, a indexação da expressão com a página correta.
% Veja os exemplos:

%\index{Convergentes de frações contínuas simples} %
%\index{Frações contínuas simples!convergentes}%

%\index{Wallis!teorema de} %
% Veja exemplo no arquivo cap1.tex

%%%%%%%%%%%%%%%%%%%%%%%%%%%%%%%%%%%%%%%%%%%%%%%%%%%%%%%%%%%%%%%%%%%%%%%%%%%%%%%%%%%%%%%%%%%%%%%%%

\end{document}
